\documentclass[11pt]{article}
\usepackage[margin=1in]{geometry}
\usepackage[T1]{fontenc}
\usepackage{lmodern}
\usepackage{longtable}
\usepackage{array}
\usepackage{enumitem}
\usepackage{hyperref}
\hypersetup{colorlinks=true,linkcolor=blue,urlcolor=blue}
\setlist[itemize]{leftmargin=1.4em}
\setlist[enumerate]{leftmargin=1.7em}
\newcommand{\code}[1]{\texttt{#1}}
\newcommand{\fpath}[1]{\texttt{#1}}
\newcommand{\st}{\ensuremath{\mathrm{STV}}}
\newcommand{\ec}{\ensuremath{\mathrm{EC}}}

\title{PeTTaChainer Codebase Assessment:\\Strengths, Weaknesses, and Repair Directions}
\author{Prepared for Ben Goertzel and the PeTTaChainer author}
\date{2026-07-04}

\begin{document}
\maketitle

\begin{abstract}
This note is an author-facing assessment of the checked-out PeTTaChainer
codebase used in the petta-memory project.  It is intentionally written as an
ASCII LaTeX source file so it can be edited, diffed, mailed, and compiled with
ordinary LaTeX tooling.  The assessment is based on source inspection of
\fpath{projects/petta-memory/repos/PeTTaChainer} at commit
\code{e4db5cad60a39c0d3f81a07296d606af6de4d76d}, on the sibling PeTTa checkout
at commit \code{d8d46920269ced70cd6236a5182d4d2409c1e12b}, and on the local
petta-memory integration probes through commit \code{cda5cbe}.  The short
version is: PeTTaChainer has an unusually valuable semantic design for
contextual, paraconsistent PLN, but the current runtime path has a serious
materialization/compilation bottleneck that makes it risky as the near-term
production chainer for petta-memory until the add path is instrumented and
hardened.
\end{abstract}

\tableofcontents

\section{Executive summary}

PeTTaChainer is not just a toy wrapper around MeTTa.  It contains several ideas
that are directly aligned with the requirements of an experiential memory
system: proof atoms, explicit truth values, evidence-count packets, generated
local contexts, context-sensitive projection to PLN truth values, forward and
backward chaining, proof audit machinery, and benchmark/verifier scripts.  In
particular, the context layer is a good match for pi-PLN style reasoning in
which evidence is not forced into a single global probability space.

The main concern is operational rather than conceptual.  In local petta-memory
probes, small statements pass syntax checks and the PeTTa/SWI/Janus runtime can
initialize, but the public add path reaches \code{compileadd}, whose first
binding \code{materialize-stmt-lambdas} can time out even on lambda-free tiny
proof atoms.  The latest narrowed failing shape is a four-field wrapper with
adjacent arity-two nested sibling payloads, for example a proof-like envelope
containing a two-argument nested statement type beside a two-argument truth-value
payload.  This appears to be a generic evaluator/materializer traversal issue,
not merely a petta-memory token issue or an STV-specific semantic error.

My recommendation is therefore a two-track strategy:
\begin{enumerate}
  \item Preserve PeTTaChainer as the semantic reference implementation for
  pi-PLN context/evidence ideas.
  \item In parallel, use a more operationally stable PLN chainer for immediate
  petta-memory experiments, and add PeTTaChainer-compatible evidence/context
  semantics around it.
\end{enumerate}

For the PeTTaChainer author, the highest-leverage repair is to isolate and
benchmark the \code{materialize-stmt-lambdas} and \code{mm2compile} stages on a
minimal failing term, then add a supported precompiled or bulk-add API so callers
can load already-normalized facts without passing every insertion through the
full dynamic compilation path.

\section{Scope and evidence}

\subsection{Repository and commit scope}

This assessment covers the local checkout:
\begin{itemize}
  \item \fpath{projects/petta-memory/repos/PeTTaChainer}
  \item Commit: \code{e4db5cad60a39c0d3f81a07296d606af6de4d76d}
  \item Commit message: \code{Implement piPLN: paraconsistent contextual reasoning over PeTTaChainer}
  \item Sibling PeTTa dependency: \fpath{projects/petta-memory/repos/PeTTa}
  \item Sibling PeTTa commit: \code{d8d46920269ced70cd6236a5182d4d2409c1e12b}
\end{itemize}

The local checkout is clean at the time of this report.  The PeTTaChainer tree
contains about 98 Python, MeTTa, Markdown, and Prolog source/documentation files
under the \fpath{pettachainer/} package tree, plus examples and tests.

\subsection{Integration evidence from petta-memory}

The petta-memory integration branch has accumulated a bounded evidence trail.
The important observations are:
\begin{itemize}
  \item SWI-Prolog 9.3.36 with Janus was installed in user space and verified.
  \item \code{PeTTa().process\_metta\_string('!(+ 2 3)')} returned \code{5}.
  \item \code{pettachainer.check\_stmt} accepted a simple statement of the form
  \code{(: bird\_robin (Bird robin) (STV 1.0 0.99))}.
  \item \code{PeTTaChainer} construction is relatively fast in the local probe,
  around half a second after setup.
  \item \code{check\_stmt} succeeds for generated petta-memory STV statements.
  \item Public add/query routes through \code{compileadd} or
  \code{compileadd-mine}; no public precompiled-add/cache API was found by
  source inspection.
  \item Bounded probes repeatedly time out inside the add-stage materialization
  path before reliable \code{mm2compile}, full \code{compileadd}, query, or
  contextual query gates can be claimed for petta-memory inputs.
\end{itemize}

The latest local petta-memory gate, at commit \code{cda5cbe}, narrows the
materializer blocker to adjacent arity-two nested sibling payloads in a
four-field wrapper.  In the petta-memory repository, \code{PYTHONPATH=src
python3 -m unittest discover -s tests -v} passes 116 tests, including the
non-live gates that preserve the boundary between checked handoff evidence and
claimed inferred beliefs.

\section{Architectural overview}

\subsection{Python surface}

The public Python package is small and useful.  The top-level interface exports:
\begin{itemize}
  \item \code{PeTTaChainer}
  \item \code{ContextualQueryResult}
  \item \code{check\_stmt}
  \item \code{check\_query}
  \item \code{get\_language\_spec}
\end{itemize}

The main class in \fpath{pettachainer/pettachainer.py} wraps a PeTTa runtime,
loads \fpath{pettachainer/metta/petta\_chainer.metta}, creates a unique KB name,
loads a logic config, and exposes methods for adding atoms, querying, forward
chaining, and contextual querying.  The wrapper also stores added atoms so a
contextual query can temporarily inject them into \code{\&self} and call the
context projection heads.

\subsection{MeTTa runtime layout}

The runtime organization is much better than a flat script pile:
\begin{itemize}
  \item \fpath{pettachainer/metta/petta\_chainer.metta}: import root and public
  MeTTa heads such as \code{compileadd}.
  \item \fpath{pettachainer/metta/chainer/}: compilation, truth-value formulas,
  distribution formulas, forward and backward chaining, proof-store helpers,
  query runtime, mining, and proof-structure audit.
  \item \fpath{pettachainer/metta/context/}: context extraction, context
  generation, pure-MeTTa and Prolog beam search, and inference-control modules.
  \item \fpath{pettachainer/metta/logic\_configs/}: selectable logic configs,
  currently including PLN and predicate-logic variants.
  \item \fpath{pettachainer/metta/piPLN\_paper\_explained/}: runnable examples
  and paper-translation cells.
\end{itemize}

\subsection{Core statement and query model}

The language spec uses proof atoms:
\begin{verbatim}
(: proof-id type tv)
\end{verbatim}

Statements are inserted with:
\begin{verbatim}
!(compileadd kb (: proof-id type tv))
\end{verbatim}

Queries use:
\begin{verbatim}
!(query steps kb (: $prf type $tv))
\end{verbatim}

The explicit \code{proof-id} slot is a strength because it gives proof objects a
place to live.  The \code{tv} slot currently supports ordinary \code{STV} values
and several value-distribution forms such as \code{PointMass}, \code{NatDist},
\code{FloatDist}, and \code{ParticleDist}.  The distinction between truth
uncertainty and value uncertainty is documented clearly in the language spec.

\section{Major strengths}

\subsection{The pi-PLN semantic idea is the right one for this use case}

The most important strength is conceptual.  PeTTaChainer represents evidence as
context-indexed positive and negative counts, \code{(EC pos neg)}, and only
projects these counts to \code{STV} values when a local query context demands it.
This is a strong fit for experiential memory, where contradictory observations
should usually remain visible rather than be prematurely averaged away.

The generated-context layer is particularly relevant.  The code enumerates
candidate guards from evidence-packet features, splits evidence inside/outside a
guard, scores the split by conflict reduction plus a query-focus bonus and a
feature penalty, then projects selected evidence to an \code{STV}.  This directly
supports exception handling such as birds/penguins without assuming a single
universal distribution.

\subsection{The code is runnable as a package, not only as a notebook}

The repository has been organized as an installable Python package with bundled
MeTTa runtime data.  The \fpath{pyproject.toml} declares the package, package
data is selected through \fpath{MANIFEST.in}, and the README gives installation
instructions.  This matters because many experimental reasoning repositories
never cross the boundary from demo scripts to importable libraries.

The Python API is also small enough to be intelligible.  A user can construct a
handler, add atoms, query, forward-chain, or ask for a contextual query without
learning every internal MeTTa helper first.

\subsection{The repository includes author-facing language documentation}

The files \fpath{pettachainer/LANGUAGE\_SPEC.md} and
\fpath{pettachainer/LLM\_RULE\_SPEC.md} are useful.  They document the core proof
atom syntax, rule templates, query patterns, truth/value uncertainty distinction,
distribution truth values, and contextual-query idea.  This is a practical
strength for collaboration with LLM agents, because it constrains generated
rules to a smaller and more checkable language.

\subsection{There is a real validation seam}

\fpath{pettachainer/pln\_validator.py} provides \code{check\_stmt} and
\code{check\_query}.  They are not deep semantic validators, but they catch
important shape errors: statement proof IDs must not be variables, query proof
slots should be variables, and truth-value forms must be in the supported set.
The local petta-memory integration used this seam successfully to validate
exported STV proof statements before treating them as PLN-ready handoff inputs.

\subsection{The chainer has both forward and backward capabilities}

The API and tests exercise forward chaining, backward queries, seeded forward
chaining from selected facts, and logic-config differences.  The query runtime
compiles a query into a goal plus temporary adds, injects those adds, runs the
chainer, externalizes proof structure, and cleans up temporary query refs.  The
cleanup path is important: without it, query compilation can leave confusing
state behind.

\subsection{Proof and audit concerns are unusually prominent}

The repository contains proof-structure audit modules, showcase verification,
forensic packets, witness certificates, red-team mutation tests, and replay
checks.  Even if some of this is heavier than petta-memory needs immediately,
the orientation is good: the codebase treats proof claims as artifacts that can
be replayed and attacked, not as opaque strings printed by a demo.

\subsection{Benchmark coverage is broader than usual}

The benchmark suite is not limited to wall-clock timing.  It includes smart
dispatch, generated context control, complementary evidence, forward/backward
comparisons, materialized backward queries, particle-vs-exact distributions,
incident response, showcase verification, and red-team forgeries.  This gives
the author several places to attach regression tests after fixing the current
materialization bottleneck.

\subsection{The codebase already separates semantic layers}

There is a meaningful separation between:
\begin{itemize}
  \item syntax and validation,
  \item compilation and internal proof representation,
  \item forward/backward chaining,
  \item truth-value and distribution formulas,
  \item context generation and context projection,
  \item benchmarks and forensic verification.
\end{itemize}

This separation is imperfect, but it is far better than putting all semantics in
one long interpreter file.  It should make targeted repair possible.

\section{Major weaknesses and risks}

\subsection{The public add path is brittle}

The public insertion methods call \code{compileadd} or \code{compileadd-mine}.
The current \code{compileadd} path is:
\begin{verbatim}
(= (compileadd $kb $stmt)
   (let*
      (($stmt1 (materialize-stmt-lambdas $stmt))
        ($atoms (collapse (mm2compile $kb $stmt1)))
        ($iatoms (list_to_set (map-flat internalize-proof-structure $atoms)))
        ($eatoms (map-flat externalize-proof-structure $iatoms))
        ($_index (index-source-implication $kb $stmt1))
        ($_ (collapse (foreach-flat add-to-kb $iatoms)))
        ($_ (maybe-process-on-add $kb $stmt1)))
      (superpose $eatoms)))
\end{verbatim}

This is a large amount of work for every atom insertion.  For small interactive
examples it may be fine.  For memory-system ingestion, batch loading, and
agent-generated evidence packets, it is risky unless each stage has cheap and
well-bounded behavior.

\subsection{The materializer can time out before real compilation begins}

The most serious observed issue is that \code{materialize-stmt-lambdas} can be
slow or nonterminating on lambda-free expressions.  Source inspection says the
materializer should only evaluate a term when the expression head is \code{|->};
otherwise it recursively rebuilds the term:
\begin{verbatim}
(= (materialize-stmt-lambdas $term)
   (if (is-var $term) $term
      (if (is-expr $term)
         (if (== (car-atom $term) |->) (eval $term)
            (cons (materialize-stmt-lambdas (car-atom $term))
               (map-flat materialize-stmt-lambdas (cdr-atom $term))))
         $term)))
\end{verbatim}

However, local probes show timeouts on lambda-free proof atoms.  The failure has
been narrowed beyond the original tokens.  Empty and one-argument nested types
pass; generic two-argument nested expressions inside certain full proof/envelope
contexts can block.  The latest failing boundary is adjacent arity-two nested
siblings in a four-field wrapper.

This means that the practical blocker is probably in PeTTa/MeTTa evaluation of
this recursive traversal pattern, in interaction with list/expression structure,
not in user lambda execution.

\subsection{No supported precompiled-add or bulk-load API was found}

Source inspection found public add methods routing through \code{compileadd} or
\code{compileadd-mine}.  No public API was found for adding already-compiled or
already-normalized facts while preserving the intended indexes/proof semantics.
This leaves integrators with bad choices:
\begin{itemize}
  \item use \code{compileadd} and hit the materialization bottleneck;
  \item directly assert low-level facts and risk bypassing important indexes;
  \item keep a non-live handoff cache outside PeTTaChainer, which is safe but
  does not exercise PeTTaChainer inference.
\end{itemize}

A supported precompiled/bulk-add interface would make the codebase much more
usable for systems that already have validated proof atoms.

\subsection{Runtime dependency setup is heavy}

The runtime requires PeTTa, SWI-Prolog 9.3.x, Janus, and in this environment
careful environment variables such as \code{SWIPL\_HOME}, \code{SWI\_HOME\_DIR},
\code{LD\_LIBRARY\_PATH}, and sometimes \code{LD\_PRELOAD}.  This is acceptable
for a research prototype but fragile for collaborator onboarding, CI, and
agent-run experiments.  The repository documents the requirement, but the cost
remains real.

\subsection{Timeout protection is incomplete around adds}

The Python \code{query} method has a multiprocessing timeout path when possible,
but \code{add\_atom} and \code{add\_atoms\_no\_check} call PeTTa directly.  If
\code{compileadd} stalls, the caller can hang unless the caller independently
runs it in a subprocess.  The petta-memory probes had to add subprocess isolation
outside the codebase for this reason.

\subsection{The global runtime load flag may hide lifecycle problems}

\code{LOADEDLIB} and \code{LOADED\_LOCK} avoid repeated library loads in a Python
process.  This is understandable, but it means that runtime state is partly
global, while each handler also creates its own KB name and PeTTa instance.  The
interaction between global MeTTa library load state, SWI/Janus state, spawned
query workers, and repeated tests deserves careful stress testing.  Bugs here
could appear only after long-running agent sessions.

\subsection{Contextual query temporarily injects atoms into \&self}

The contextual query path copies stored atoms into \code{\&self}, calls the
context projection machinery, and removes them in a \code{finally} block.  The
\code{finally} is good.  But temporary mutation of a shared space is inherently
more fragile than querying a separate explicit evidence space.  If duplicate
atoms, removal failures, concurrent calls, or interrupted runtimes occur, context
projection may see stale evidence.

\subsection{The semantic surface is larger than the validation surface}

The validator checks top-level statement/query shape and supported TV forms.  It
does not fully validate rule schemas, evidence packet provenance, variable
safety, context features, distribution payloads, proof IDs, or whether generated
terms will be safe in downstream compiler stages.  That is acceptable for a thin
syntax guard, but callers should not mistake it for a semantic type checker.

\subsection{Benchmarks may not isolate the failing micro-stage}

The benchmark suite is broad, but the current petta-memory blocker needs very
small stage-isolated tests: one term, one materializer, bounded output, bounded
timeout.  The existing showcases prove many positive paths, but they do not by
themselves explain why a tiny lambda-free proof atom can time out inside
\code{materialize-stmt-lambdas}.  A small failing regression test should be added
near the materializer itself.

\section{Detailed observations by subsystem}

\subsection{Packaging and installation}

\textbf{Strengths.}  The package metadata is clear, the MIT license is present,
the package name and project URLs are defined, and the README explains that
PeTTa is a sibling dependency rather than a PyPI dependency.  Bundling the MeTTa
runtime as package data is the right move for a Python library.

\textbf{Weaknesses.}  The dependency chain remains nontrivial, especially the
SWI/Janus build.  Continuous integration should explicitly cover clean-venv
installation and a tiny smoke under the documented SWI version.  If CI cannot
run the full stack, the README should say exactly which smoke was last run and
on what platform.

\subsection{Python API}

\textbf{Strengths.}  The API is concise.  \code{ContextualQueryResult.answers}
puts the generated projection first, then ordinary proofs.  Query timeouts use a
fresh spawned process when possible, which avoids inheriting live SWI/Janus state
into a forked child.

\textbf{Weaknesses.}  Adds lack the same timeout isolation as queries.  In
addition, \code{add\_atoms\_no\_check} batches all additions into a single
\code{superpose}; if one item triggers the materializer bottleneck, the entire
batch becomes hard to diagnose.  A safer API would provide per-atom staged
results and optional subprocess isolation.

\subsection{Compilation path}

\textbf{Strengths.}  The compilation pipeline has named stages: materialization,
MeTTa-to-compiled conversion, proof internalization/externalization, indexing,
KB insertion, and on-add processing.  This is good because it gives us stage
names for profiling and repair.

\textbf{Weaknesses.}  The stages are currently fused into a single public
operation.  There is no obvious public way to run them individually with trace,
step budget, timeout, or return structured diagnostics.  For research tooling,
\code{compileadd-explain} or \code{compileadd-stages} would be very valuable.

\subsection{Query runtime}

\textbf{Strengths.}  The query runtime injects temporary compiled adds, runs the
chainer, externalizes proofs, pretty-prints, drops empty results, and clears
temporary query additions before and after.  This is the right overall pattern.

\textbf{Weaknesses.}  The same compiled-query path depends on \code{mm2compileQuery}
and the broader compiler stack.  If the compiler is brittle for certain nested
shapes, query reliability will be shape-sensitive too.  The temporary ref cleanup
should be stress-tested under exceptions and nested queries.

\subsection{Context generation}

\textbf{Strengths.}  The context layer is the strongest part of the codebase for
AGI research use.  It explicitly encodes evidence packets, guard enumeration,
conflict measurement, split scoring, focus bonuses, feature penalties, evidence
selection, STV projection, proof terms, and minimality reports.  It also has both
pure-MeTTa and Prolog beam variants.

\textbf{Weaknesses.}  The scoring constants and guard-enumeration bounds are
research choices that need more explicit calibration.  The current pure
enumerator tops out at depth 3, while beam variants can go deeper.  That is fine,
but the user-facing API should make the selected search regime and its limits
obvious in every result.

\subsection{Truth-value and evidence semantics}

\textbf{Strengths.}  The code keeps \code{EC} counts separate from \code{STV}
projection, and the language spec distinguishes truth uncertainty from value
uncertainty.  This is exactly the kind of distinction that prevents subtle PLN
modeling errors.

\textbf{Weaknesses.}  The mapping from raw evidence/provenance to \code{EC}
counts remains the responsibility of the caller or data source.  That is not a
bug, but a documentation and API risk.  If users infer \code{EC} counts from
\code{STV} confidence without a principled rule, the pi-PLN semantics will be
polluted.

\subsection{Tests and verification}

\textbf{Strengths.}  The Python tests cover forward/backward behavior, contextual
queries, validators, context showcases, complementary evidence, incident
response, smart dispatch, and forensic verification.  This is unusually rich
for a research prototype.

\textbf{Weaknesses.}  The test suite contains heavy showcase tests that may be
too slow for every edit.  It needs a small, fast ``compiler microtests'' layer
that isolates specific recursive compiler/materializer shapes.  The current
petta-memory failing shape should become one such microtest.

\section{Minimal failing shape to hand to the author}

The following is not claimed as the only failing shape, but it captures the
boundary found by petta-memory instrumentation.

Shapes that tend to pass quickly:
\begin{itemize}
  \item a nested type expression alone;
  \item a full wrapper with an atom type;
  \item a full wrapper with an empty nested type expression;
  \item a full wrapper with a one-argument nested type expression;
  \item generic four-field wrappers where only one relevant nested sibling has
  arity two.
\end{itemize}

Shapes that can block under bounded probes:
\begin{itemize}
  \item a four-field proof/envelope-like expression containing a two-argument
  nested type beside an \code{STV}-like two-argument payload;
  \item the generalized form where adjacent sibling payloads are both nested
  expressions of arity two.
\end{itemize}

A representative term family is:
\begin{verbatim}
(ProofEnvelope PayloadA (Requires TypeArgSentinel0 TypeArgSentinel1)
              (RightPayload 1.0 1.0))
\end{verbatim}

and the petta-memory-relevant proof atom is:
\begin{verbatim}
(: b-profile-000 (Requires MemoryTarget0 PLNReadyViews) (STV 1.0 1.0))
\end{verbatim}

The author should first test \code{materialize-stmt-lambdas} directly on these
terms, outside the full \code{compileadd} path, with timing and recursion trace.
If direct materialization is slow, the bug is before \code{mm2compile}.  If it is
fast in isolation but slow inside \code{compileadd}, then the interaction with
\code{let*}, \code{collapse}, \code{map-flat}, or runtime evaluation context is
suspect.

\section{Recommended fixes and improvements}

\subsection{Priority 1: stage-isolated compiler diagnostics}

Add a diagnostic head, or a Python helper, that runs the public add pipeline one
stage at a time:
\begin{enumerate}
  \item validate/evaluate statement;
  \item materialize lambdas;
  \item run \code{mm2compile};
  \item internalize proof structure;
  \item externalize proof structure;
  \item index source implication;
  \item add internalized atoms to the KB;
  \item run on-add processing.
\end{enumerate}

Each stage should return structured information: status, elapsed time, output
size, abbreviated output hash, and error/timeout.  This can be implemented in
Python with subprocess boundaries if MeTTa-level timeouts are unavailable.

\subsection{Priority 2: repair or replace recursive materialization}

The current recursive definition of \code{materialize-stmt-lambdas} is elegant
but appears operationally dangerous.  Possible repairs:
\begin{itemize}
  \item Add a fast path: if a statement contains no \code{|->}, return it
  unchanged without recursively rebuilding the whole expression.
  \item Add a structural scanner that detects lambda forms without invoking
  expensive expression reconstruction.
  \item Replace \code{map-flat} traversal with a more predictable list walk if
  \code{map-flat} interacts badly with nested expression siblings.
  \item Add a depth/node budget and return an explicit error rather than hanging.
\end{itemize}

For petta-memory's current statements, a no-lambda fast path should make
\code{materialize-stmt-lambdas} an identity operation.

\subsection{Priority 3: add a supported precompiled or bulk-add API}

Many integrators will have already-normalized facts and rules.  A supported API
should let them load those facts while preserving whatever indexes and proof
structures PeTTaChainer requires.  The API could be named something like:
\begin{verbatim}
compileadd-checked
compileadd-precompiled
bulk-compileadd
kb-add-compiled-proof-atoms
\end{verbatim}

The key is not the name.  The key is that callers should not be forced to use
private low-level assertions that might bypass indexes.  The API should state
which invariants it maintains.

\subsection{Priority 4: add timeout-safe add methods to the Python API}

Mirror the query timeout design for \code{add\_atom} and \code{add\_atoms}.  A
safe method could return a structured object instead of raising raw exceptions:
\begin{verbatim}
AddResult(status, stage, elapsed_s, output, error, diagnostics)
\end{verbatim}

This would make failures much easier for agent systems to report.

\subsection{Priority 5: strengthen semantic validation}

Keep \code{check\_stmt} and \code{check\_query}, but add optional deeper checks:
\begin{itemize}
  \item proof-id symbol validity;
  \item rule variable safety;
  \item supported implication/premise/conclusion schemas;
  \item evidence-packet shape and nonnegative counts;
  \item context feature-list shape;
  \item distribution payload well-formedness;
  \item compiler-safety shape checks for known problematic nested forms.
\end{itemize}

This can be staged.  A thin validator is useful, but a research integration path
needs a ``strict mode'' validator too.

\subsection{Priority 6: expose context search limits in outputs}

Generated context results should always record which search method was used
(pure enumeration, pure beam, Prolog beam), maximum depth, beam width, candidate
count if available, selected guard, runner-up guard, score margin, and evidence
cap/prior.  This will make context projections more scientifically auditable.

\section{Recommended petta-memory stance}

For petta-memory, I recommend the following boundary until PeTTaChainer add/query
is fixed:
\begin{enumerate}
  \item Continue exporting promotion-eligible beliefs as checked STV proof
  statements and explicit-count EvidencePackets.
  \item Keep the handoff cache labelled as input evidence, not inferred belief.
  \item Do not claim PeTTaChainer inference over petta-memory until a bounded
  \code{compileadd} or supported equivalent passes on the tiny fixture.
  \item Use PeTTaChainer semantics as a design reference for pi-PLN evidence and
  context semantics.
  \item In parallel, prototype pi-PLN evidence/context wrappers around a more
  functional PLN chainer, currently \fpath{patham9/PLN}, for near-term experiments.
\end{enumerate}

This is not a rejection of PeTTaChainer.  It is a way to avoid letting one
runtime bottleneck block the conceptual and integration progress.

\section{Suggested regression tests for PeTTaChainer}

\subsection{Fast materializer identity tests}

Add direct tests for \code{materialize-stmt-lambdas} on lambda-free terms:
\begin{verbatim}
(Requires MemoryTarget0 PLNReadyViews)
(STV 1.0 1.0)
(: b-profile-000 (Requires MemoryTarget0 PLNReadyViews) (STV 1.0 1.0))
(ProofEnvelope PayloadA (Requires TypeArgSentinel0 TypeArgSentinel1)
              (RightPayload 1.0 1.0))
\end{verbatim}

Expected result: identity output under a small time bound.

\subsection{Stage timing tests}

For a tiny fact and a tiny one-premise implication, record time for each
\code{compileadd} stage.  Fail if any stage exceeds a conservative bound in CI,
or at least mark it as a performance regression.

\subsection{Batch-add diagnostics}

Add a test that loads a list of atoms and returns per-atom status.  The first
failing atom should be identifiable without bisecting by hand.

\subsection{Context projection reproducibility}

For the bird/penguin example, record not only the final low-STV projection but
also selected guard, support packets, score trace, and minimality report.  Verify
that irrelevant noise does not alter the selected guard unless the score margin
is genuinely near zero.

\section{Strengths and weaknesses table}

\begin{longtable}{p{0.23\linewidth}p{0.34\linewidth}p{0.34\linewidth}}
\textbf{Area} & \textbf{Strength} & \textbf{Weakness or risk} \\
\hline
Semantic model & Context-indexed EC evidence, local STV projection, generated
contexts, paraconsistent conflict handling. & Needs careful documentation of how
raw evidence becomes EC counts; users may misuse STV confidence as EC evidence.
\\
Python API & Small importable interface; query timeout support; contextual
query result object. & Add path has no equivalent timeout isolation; temporary
\&self injection is fragile under concurrency/interruption. \\
Runtime layout & Clear split into chainer, context, configs, tests, benchmarks,
and paper examples. & Some critical public behavior still depends on large fused
MeTTa heads. \\
Compiler path & Named stages make repair possible. & \code{compileadd} fuses
all stages; \code{materialize-stmt-lambdas} can block before compilation.
\\
Validation & Useful top-level statement/query shape checks. & Not a strict
semantic type checker; does not catch all rule/evidence/context risks. \\
Tests & Broad showcase, context, audit, verifier, and red-team coverage. & Needs
small microtests for compiler/materializer shape regressions. \\
Packaging & Installable Python package with bundled MeTTa runtime data. & Heavy
SWI/Janus/PeTTa dependency setup; CI/onboarding remain fragile. \\
Benchmarks & Covers dispatch, contexts, evidence merge, incident replay, and
forensic packets. & Current failing petta-memory case needs more stage-isolated
benchmarks than broad showcases. \\
\end{longtable}

\section{Conclusion}

PeTTaChainer is semantically promising and unusually close to what an
OmegaClaw/petta-memory style experiential reasoner wants: context-sensitive,
paraconsistent, proof-aware PLN over a MeTTa/PeTTa substrate.  Its main weakness
is that the operational add/compile path is not yet robust enough for dependable
agent memory integration.  The codebase should not be discarded; it should be
stabilized around the materialization/compiler boundary and given a safe staged
or precompiled load API.

Until then, petta-memory should preserve PeTTaChainer-compatible handoff
artifacts and use the codebase as the semantic reference, while developing the
near-term pi-PLN chainer path on a more functional substrate.  This keeps the
research direction intact without overstating current runtime maturity.

\appendix
\section{Key source paths}

\begin{itemize}
  \item \fpath{pettachainer/pettachainer.py}: Python handler and API.
  \item \fpath{pettachainer/pln\_validator.py}: statement/query shape checks.
  \item \fpath{pettachainer/LANGUAGE\_SPEC.md}: user-facing language spec.
  \item \fpath{pettachainer/LLM\_RULE\_SPEC.md}: LLM-oriented rule construction spec.
  \item \fpath{pettachainer/metta/petta\_chainer.metta}: MeTTa entry point,
  \code{materialize-stmt-lambdas}, \code{compileadd}.
  \item \fpath{pettachainer/metta/chainer/compile.metta}: compiler logic.
  \item \fpath{pettachainer/metta/chainer/compiled\_query\_runtime.metta}: query runtime.
  \item \fpath{pettachainer/metta/context/context\_generation.metta}: generated
  contexts and EC-to-STV projection.
  \item \fpath{pettachainer/metta/context/context\_from\_kb.metta}: context
  evidence extraction from a KB space.
  \item \fpath{pettachainer/metta/context/context\_inference\_control.metta} and
  \fpath{context\_inference\_control\_beam.metta}: context-driven branch control.
  \item \fpath{pettachainer/metta/piPLN\_paper\_explained/}: runnable reading
  guide and demos.
  \item \fpath{pettachainer/benchmarks/}: benchmark and verification scripts.
  \item \fpath{tests/}: Python regression tests.
\end{itemize}

\section{Petta-memory artifacts informing this report}

The following local artifacts and decisions informed the assessment.  They are
not all required reading for the PeTTaChainer author, but they record the
integration evidence trail.

\begin{itemize}
  \item \fpath{projects/petta-memory/DECISIONS.md}: PeTTaChainer integration,
  compileadd bottleneck, static-import side path, materializer gates, and pivot
  to \fpath{patham9/PLN} as a parallel functional chainer base.
  \item \fpath{projects/petta-memory/TASKS.md}: current tasks and latest status.
  \item \fpath{projects/petta-memory/repos/petta-memory/artifacts/}:
  structured JSON outputs from bounded profiling and materializer probes.
  \item Latest relevant materializer artifact:
  \fpath{pettachainer\_materialize\_four\_field\_adjacent\_nested\_arity\_gate\_2026-07-05T0400Z.json}.
\end{itemize}

\end{document}
